% =====================================================================
% Actividad práctica - Ingeniería de Requisitos (ISR-401) -- UTEQ
% Adaptación al formato institucional del documento de referencia.
% El contenido de las actividades corresponde únicamente al trabajo realizado.
% =====================================================================
\documentclass[11pt,a4paper,oneside]{article}

\usepackage[utf8]{inputenc}
\usepackage[T1]{fontenc}
\usepackage[spanish,es-nodecimaldot]{babel}
\renewcommand{\tablename}{Tabla}
\renewcommand{\figurename}{Figura}
\usepackage[scaled=0.92]{helvet}
\usepackage{textcomp}

\usepackage[a4paper,top=2.4cm,bottom=2.3cm,left=2.5cm,right=2.3cm,headheight=15pt]{geometry}

\usepackage{amsmath,amssymb}
\usepackage{graphicx}
\usepackage[table,dvipsnames]{xcolor}
\usepackage{array}
\usepackage{tabularx}
\usepackage{multirow}
\usepackage{colortbl}
\usepackage{booktabs}
\usepackage{enumitem}
\usepackage[protrusion=true,expansion=false]{microtype}
\usepackage{parskip}
\usepackage{titlesec}
\usepackage{fancyhdr}
\usepackage{caption}
\usepackage{pdflscape}
\usepackage[most]{tcolorbox}
\usepackage{longtable}
\usepackage{float}
\usepackage{tikz}
\usetikzlibrary{positioning,arrows.meta,shapes.geometric,shapes.multipart,calc}

% Paleta institucional
\definecolor{darkblue}{HTML}{1A3A5C}
\definecolor{midblue}{HTML}{2E6DA4}
\definecolor{gold}{HTML}{C8A951}
\definecolor{lightgray}{HTML}{F2F2F2}
\definecolor{rowgray}{HTML}{EEEEEE}
\definecolor{softblue}{HTML}{EAF1F8}
\definecolor{softred}{HTML}{F7E7E7}
\definecolor{darkred}{HTML}{9C2A2A}
\definecolor{softgreen}{HTML}{E7F1E7}

\linespread{1.05}
\setlength{\parskip}{5pt}
\setlength{\parindent}{0pt}
\renewcommand{\arraystretch}{1.25}
\newcolumntype{L}[1]{>{\raggedright\arraybackslash}p{#1}}
\newcolumntype{C}[1]{>{\centering\arraybackslash}p{#1}}
\newcolumntype{R}[1]{>{\raggedleft\arraybackslash}p{#1}}

\titleformat{\section}{\sffamily\large\bfseries\color{darkblue}}{\thesection}{0.6em}{}
[\vspace{-5pt}{\color{gold!85}\rule{\linewidth}{1pt}}]
\titleformat{\subsection}{\sffamily\normalsize\bfseries\color{midblue}}{\thesubsection}{0.5em}{}
\titleformat{\subsubsection}{\sffamily\normalsize\bfseries\color{midblue}}{\thesubsubsection}{0.5em}{}

\titlespacing*{\section}{0pt}{14pt}{6pt}
\titlespacing*{\subsection}{0pt}{10pt}{3pt}

\pagestyle{fancy}
\fancyhf{}
\renewcommand{\headrulewidth}{0.4pt}
\fancyhead[L]{\footnotesize\sffamily\color{darkblue}Actividad Práctica\, \textbullet\, Unidad IV}
\fancyhead[R]{\footnotesize\sffamily\color{darkblue}Ingeniería de Requisitos\, \textbullet\, ISR-401}
\fancyfoot[C]{\footnotesize\sffamily\color{darkblue}\thepage}

\captionsetup{font=small,labelfont={bf,sf,color=darkblue},labelsep=period,skip=5pt}

\newtcolorbox{concepto}[1]{
  colback=softblue,colframe=darkblue,boxrule=0.6pt,arc=2.5pt,
  left=7pt,right=7pt,top=5pt,bottom=5pt,
  fonttitle=\bfseries\sffamily\color{white},
  coltitle=white,colbacktitle=darkblue,title={#1},breakable
}

\newcommand{\hd}[1]{\cellcolor{darkblue}\textbf{\textcolor{white}{\small\sffamily #1}}}
\newcommand{\hm}[1]{\cellcolor{midblue}\textbf{\textcolor{white}{\small\sffamily #1}}}
\newcommand{\hg}[1]{\cellcolor{lightgray}\textbf{\color{darkblue}\small\sffamily #1}}

\begin{document}
\sffamily

% Portada breve
\begin{center}
  {\Large\bfseries\color{darkblue} UNIVERSIDAD TÉCNICA ESTATAL DE QUEVEDO}\\[2pt]
  {\normalsize\color{midblue} Facultad de Ciencias de la Ingeniería \textbullet\ Carrera de Ingeniería de Software}\\[6pt]
  {\color{gold}\rule{0.9\linewidth}{1.2pt}}\\[6pt]
  {\large\bfseries\color{darkblue} ACTIVIDAD PRÁCTICA --- UNIDAD IV}\\[2pt]
  {\normalsize\color{darkblue} Validación, Gestión, Herramientas y Estándares de Requisitos}\\[2pt]
  {\small\color{black} Asignatura: Ingeniería de Requisitos (ISR-401)}
\end{center}

\vspace{2pt}
\begin{center}\small
\begin{tabularx}{\linewidth}{@{}L{2.6cm} X L{2.6cm} X@{}}
  \hg{Estudiante} & Cedeño Ávila Winston Damián & \hg{Fecha} & 12/8/26 \\
  \hg{Nivel} & 4to semestre Software ``B'' & \hg{Modalidad} & Práctica individual \\
\end{tabularx}
\end{center}

\begin{concepto}{Contenido presentado}
El documento conserva únicamente el desarrollo realizado en la actividad transcrita, adaptándolo al formato visual y estructural del documento de referencia.
\end{concepto}

\section*{P1. Modelo de datos -- Diagrama de Clases}

% Diagrama correspondiente a P1
\begin{figure}[H]
\centering

\begin{tikzpicture}[
    class/.style={
        draw,
        rectangle,
        minimum width=5.0cm,
        text width=4.8cm,
        align=left,
        font=\small,
        inner sep=6pt
    },
    arrow/.style={-{Stealth}}
]

% ---------------- PACIENTE ----------------
\node[class] (paciente) at (0,3.5) {%
\centering\textbf{Paciente}\\[-2pt]
\noindent\rule{\linewidth}{0.4pt}\\
\raggedright
\underline{Cedula} : String\\
P.nombre : String\\
Correo : String\\[-2pt]
\noindent\rule{\linewidth}{0.4pt}\\
SolicitarCita() : Cita
};

% ---------------- CITA ----------------
\node[class] (cita) at (7,3.5) {%
\centering\textbf{Cita}\\[-2pt]
\noindent\rule{\linewidth}{0.4pt}\\
\raggedright
\underline{idCita} : int\\
FechaHora : DateTime\\
Estado : String\\[-2pt]
\noindent\rule{\linewidth}{0.4pt}\\
Cancelar() : void\\
Confirmar() : void
};

% ---------------- MEDICO ----------------
\node[class] (medico) at (0,-2.5) {%
\centering\textbf{Medico}\\[-2pt]
\noindent\rule{\linewidth}{0.4pt}\\
\raggedright
\underline{Codigo} : String\\
M.nombre : String\\
especialidad : String\\[-2pt]
\noindent\rule{\linewidth}{0.4pt}\\
Verificar Cupo(Cupo:Cupo) : bool
};

% ---------------- CUPO ----------------
\node[class] (cupo) at (7,-2.5) {%
\centering\textbf{Cupo}\\[-2pt]
\noindent\rule{\linewidth}{0.4pt}\\
\raggedright
FechaHora : DateTime\\
disponibilidad : bool\\[-2pt]
\noindent\rule{\linewidth}{0.4pt}\\
reservar() : void
};

% ---------------- RELACIONES ----------------

% Paciente -- Cita
\draw (paciente.east) -- (cita.west)
    node[pos=0.15, above] {1}
    node[pos=0.85, above] {0..*};

% Cita -- Médico
\draw (cita.south west) -- (medico.north east)
    node[pos=0.20, left] {0..*}
    node[pos=0.80, right] {1};

% Médico -- Cupo
\draw (medico.east) -- (cupo.west)
    node[pos=0.15, above] {1}
    node[pos=0.85, above] {0..*};

% Cita -- Cupo
\draw[arrow]
    (cita.east) -- ++(1.0,0)
    -- ++(0,-3.7)
    -- (cupo.east);
    
\end{tikzpicture}

\caption{Figura 1. Diagrama de clases}
\end{figure}

\section*{P2. Diagrama de actividades -- Agendar Cita}

% Diagrama correspondiente a P2
\begin{figure}[H]
\centering
\begin{tikzpicture}[
    activity/.style={
        draw,
        rectangle,
        rounded corners=2pt,
        minimum width=4.4cm,
        minimum height=1.1cm,
        align=center,
        font=\small
    },
    decision/.style={
        draw,
        diamond,
        aspect=2,
        align=center,
        font=\small,
        inner sep=8pt,
        minimum width=3.8cm
    },
    flow/.style={
        -{Stealth},
        line width=0.7pt
    },
    etiqueta/.style={
        font=\small,
        fill=white,
        inner sep=1pt
    }
]

% =========================================================
% CARRILES HORIZONTALES
% =========================================================
\draw (0,9) rectangle (19.5,0);
\draw (0,6.0) -- (19.5,6.0);
\draw (16.6,9) -- (16.6,0);
\node[font=\bfseries] at (18.05,7.5) {PACIENTE};
\node[font=\bfseries] at (18.05,3.0) {SISTEMA};

% =========================================================
% PACIENTE
% =========================================================
\node[circle, fill=black, minimum size=8pt, inner sep=0pt] (inicio) at (5,8.3) {};
\draw[flow] (inicio.south) -- (5,7.6);

\node[activity] (solicitar) at (5,7.05) {Solicitar Cita};

% Paciente -> Sistema
\draw[flow] (solicitar.south) -- (5,5.35);

% =========================================================
% SISTEMA
% =========================================================
\node[activity] (verificar) at (5,4.85) {Verificar Cupo};

% =========================================================
% DECISIÓN (alineada verticalmente con Verificar Cupo)
% =========================================================
\node[decision] (decision) at (5,3.0) {¿Cupo\\disponible?};

\draw[flow] (verificar.south) -- (decision.north);

% =========================================================
% RAMA SÍ
% =========================================================
\node[activity] (confirmar) at (2.4,1.2) {Confirmar Cita};

\draw[flow]
    (decision.west) -| node[etiqueta, pos=0.25, above] {Sí} (confirmar.north);
% =========================================================
% RAMA NO
% =========================================================
\node[activity, minimum width=5cm] (notificar) at (11.5,3.0) {Notificar no\\disponibilidad};

\draw[flow]
    (decision.east) -- node[etiqueta, midway, above] {No} (notificar.west);

% =========================================================
% PROPONER HORARIO
% =========================================================
\node[activity, minimum width=5cm] (proponer) at (11.5,1.2) {Proponer horario};

\draw[flow] (notificar.south) -- (proponer.north);

% =========================================================
% NODO FINAL (centrado entre las dos ramas)
% =========================================================
\node[circle, draw, minimum size=18pt, inner sep=0pt] (fin) at (6.95,0.4) {};
\node[circle, fill=black, minimum size=8pt, inner sep=0pt] at (fin.center) {};

\draw[flow] (confirmar.south) -- (2.4,0.4) -- (fin.west);
\draw[flow] (proponer.south) -- (11.5,0.4) -- (fin.east);

\end{tikzpicture}
\caption{Diagrama de actividades: Agendar cita}
\label{fig:diagrama_actividades_agendar_cita}
\end{figure}

\section*{P3. Maquina de estado de citas}

% Diagrama correspondiente a P3
\begin{figure}[H]
\centering
\begin{tikzpicture}[
    state/.style={
        draw,
        rectangle,
        minimum width=3.6cm,
        minimum height=0.9cm,
        align=center,
        font=\small
    },
    flow/.style={-{Stealth}, line width=0.7pt},
    etiqueta/.style={font=\small, fill=white, inner sep=2pt},
    note/.style={font=\footnotesize, align=left}
]

% =========================================================
% ESTADO INICIAL
% =========================================================
\node[circle, fill=black, minimum size=8pt, inner sep=0pt] (inicio) at (2.6,13.0) {};

% =========================================================
% SOLICITADA
% =========================================================
\node[state] (solicitada) at (2.6,11.5) {Solicitada};
\draw[flow] (inicio.south) -- (solicitada.north);

\node[note] at (7.3,11.9) {Cita creada por\\el paciente};

% =========================================================
% REGIÓN COMPUESTA "ACTIVA"
% =========================================================
\draw (0,4.6) rectangle (8.6,9.6);
\node[font=\bfseries, anchor=north west] at (0.2,9.4) {Activa};

% =========================================================
% CONFIRMADA
% =========================================================
\node[state] (confirmada) at (2.6,8.2) {Confirmada};
\draw[flow] (solicitada.south) -- node[etiqueta, midway, right] {cupo disponible} (confirmada.north);

% =========================================================
% CANCELADA (fuera de la región Activa)
% =========================================================
\node[state] (cancelada) at (12.8,8.2) {Cancelada};
\draw[flow] (confirmada.east) -- node[etiqueta, midway, above] {cancelar} (cancelada.west);

\node[note] at (12.8,9.9) {Cita cancelada\\por el paciente};

% =========================================================
% NO ASISTIDA
% =========================================================
\node[state] (noasistida) at (5.6,6.0) {No asistida};
\draw[flow] (confirmada.south) -- node[etiqueta, pos=0.35, right] {no se presenta} (noasistida.north);

\node[note] at (10.8,6.0) {Paciente no se\\presentó a la cita};

% =========================================================
% ATENDIDA
% =========================================================
\node[state] (atendida) at (2.6,1.8) {Atendida};
\draw[flow]
    (confirmada.west)
    -- ++(-1.3,0)
    |- node[etiqueta, pos=0.75, left] {atender}
    (atendida.west);

\node[note] at (-2.4,1.8) {Cita atendida\\por el médico};

% =========================================================
% NODO FINAL ÚNICO
% =========================================================
\node[circle, draw, minimum size=17pt, inner sep=0pt] (fin) at (5.6,0.3) {};
\node[circle, fill=black, minimum size=8pt, inner sep=0pt] at (fin.center) {};

\draw[flow] (atendida.south) -- (2.6,0.3) -- (fin.west);
\draw[flow] (noasistida.south) -- (5.6,0.3) -- (fin.north);
\draw[flow] (cancelada.south) -- (12.8,0.3) -- (fin.east);

\end{tikzpicture}
\caption{Máquina de estados de Cita}
\end{figure}

\section*{P4. Verificación cruzada}

\begin{longtable}{|p{0.29\textwidth}|p{0.29\textwidth}|p{0.34\textwidth}|}
\hline
\textbf{Evento P3} & \textbf{Acción P2} & \textbf{Modificación P3} \\ \hline
Solicitar Cita & Solicitar Cita & Cita. estado \\ \hline
Cupo disponible & Verificar Cupo & Cupo. disponibilidad \\ \hline
Confirmar & Confirmar Cita & Cita. estado \\ \hline
Cancelar & Cancelar cita & Cita. estado \\ \hline
Atender & Atender cita & Cita. estado \\ \hline
No se presenta & Registrar inasistencia & Cita. estado \\ \hline
\end{longtable}

\textbf{Inconsistencia encontrada:} P3 tenía ``Atender'', pero P2 no lo mostraba.

\textbf{corrección:} Se agrega la accion atender cita al proceso.

\section*{P5. Requisitos}

\subsection*{RF01 -- Agendar cita}

\begin{tabular}{|p{0.32\textwidth}|p{0.55\textwidth}|}
\hline
\textbf{Tipo} & Funcional \\ \hline
\textbf{Prioridad} & Alta \\ \hline
\textbf{Estado} & Probado \\ \hline
\textbf{Fuente} & Clínica \\ \hline
\textbf{Estabilidad} & Estable \\ \hline
\textbf{Versión} & 1.0 \\ \hline
\textbf{Verificación} & Prueba funcional \\ \hline
\end{tabular}

El sistema deberá permitir registrar una cita únicamente cuando exista un cupo disponible.

\subsection*{RF02 -- Cancelar cita}

\begin{tabular}{|p{0.32\textwidth}|p{0.55\textwidth}|}
\hline
\textbf{Tipo} & Funcional \\ \hline
\textbf{Prioridad} & Alta \\ \hline
\textbf{Estado} & Probado \\ \hline
\textbf{Fuente} & Clínica \\ \hline
\textbf{Estabilidad} & Estable \\ \hline
\textbf{Versión} & 1.0 \\ \hline
\textbf{Verificación} & Prueba funcional \\ \hline
\end{tabular}

El sistema deberá permitir cancelar una cita confirmada que aún no haya sido atendida.

\subsection*{RNF01 -- Eficiencia de desempeño}

\begin{tabular}{|p{0.32\textwidth}|p{0.55\textwidth}|}
\hline
\textbf{Prioridad} & Alta \\ \hline
\textbf{Estado} & Probado \\ \hline
\textbf{Fuente} & Clínica \\ \hline
\textbf{Estabilidad} & Estable \\ \hline
\textbf{Versión} & 1.0 \\ \hline
\textbf{Verificación} & Prueba de rendimiento \\ \hline
\end{tabular}

El sistema deberá registrar una cita en menos de 3 segundos.

\subsection*{RNF02 -- Seguridad}

\begin{tabular}{|p{0.32\textwidth}|p{0.55\textwidth}|}
\hline
\textbf{Prioridad} & Alta \\ \hline
\textbf{Estado} & Probado \\ \hline
\textbf{Fuente} & Clínica \\ \hline
\textbf{Estabilidad} & Estable \\ \hline
\textbf{Versión} & 1.0 \\ \hline
\textbf{Verificación} & Prueba de seguridad \\ \hline
\end{tabular}

El sistema deberá restringir el acceso a los datos personales únicamente a usuarios autorizados.

\section*{P6. MoSCoW}

\begin{tabular}{|p{0.20\textwidth}|p{0.20\textwidth}|p{0.48\textwidth}|}
\hline
\textbf{Req} & \textbf{Categoría} & \textbf{Justificación} \\ \hline
RF01 & Must & Es esencial para agendar citas \\ \hline
RF02 & Must & Necesario para administrar las citas \\ \hline
RNF01 & Should & Garantiza rapidez del sistema \\ \hline
RNF02 & Must & Protege los datos personales \\ \hline
\end{tabular}

\section*{P7. Validación 29148}

\begin{tabular}{|p{0.18\textwidth}|p{0.72\textwidth}|}
\hline
\textbf{Req} & \textbf{Defecto} \\ \hline
RF01 & Inicialmente no indicaba Claramente que ocurre si no existe cupo \\ \hline
RF02 & Inicialmente no especificaba que la cita debía estar sin atender \\ \hline
RNF01 & La métrica no estaba definida Inicialmente \\ \hline
RNF02 & No indicaba quién puede acceder a los datos \\ \hline
\end{tabular}

\vspace{0.3cm}

Los cuatro requisitos fueron corregidos agregando condiciones, métricas y criterios verificables.

\section*{P8. Pruebas de aceptación}

\begin{tabular}{|p{0.17\textwidth}|p{0.22\textwidth}|p{0.49\textwidth}|}
\hline
\textbf{Prueba} & \textbf{Tipo} & \textbf{Entrada} \\ \hline
PA01 & Funcional & Paciente + Cupo disponible \\ \hline
PA02 & Funcional & Cita confirmada \\ \hline
PA03 & No funcional & Solicitud de cita \\ \hline
PA04 & Seguridad & Usuario autorizado / no \\ \hline
\end{tabular}

\vspace{0.3cm}

\begin{tabular}{|p{0.22\textwidth}|p{0.50\textwidth}|p{0.18\textwidth}|}
\hline
\textbf{Prueba} & \textbf{Éxito} & \textbf{Req} \\ \hline
PA01 & Cita confirmada & RF01 \\ \hline
PA02 & Cita cancelada & RF02 \\ \hline
PA03 & Respuesta $\leq$ 3 seg & RNF01 \\ \hline
PA04 & Solo autorizado accede & RNF02 \\ \hline
\end{tabular}

\section*{P9. Matriz de trazabilidad}

\begin{tabular}{|p{0.18\textwidth}|p{0.18\textwidth}|p{0.30\textwidth}|p{0.20\textwidth}|}
\hline
\textbf{Req} & \textbf{Fuente} & \textbf{Modelo} & \textbf{Prueba} \\ \hline
RF01 & Clínica & P1, P2, P3 & PA01 \\ \hline
RF02 & Clínica & P1, P3 & PA02 \\ \hline
RNF01 & Clínica & P1 & PA03 \\ \hline
RNF02 & Clínica & P1 & PA04 \\ \hline
\end{tabular}

\vspace{0.3cm}

\textbf{Dependencia horizontal:} RF01 $\rightarrow$ RF02, porque una cita primero debe existir para poder cancelarla.

\textbf{Requisitos huérfanos:} ninguno.

\section*{P10. Gestión del cambio y línea base}

\textbf{Solicitud de cambio:} Se solicita modificar RNF01 de $\leq$ 3 seg a $\leq$ 2 seg.

\textbf{Impacto:} afecta P5, P6, P8 y P9.

\textbf{Decisión del CCB:} aprobado $\checkmark$

\textbf{Nueva versión:} RNF01 V1.1

\subsection*{Línea base}

\textbf{ERS V1.1}

\begin{tabular}{|p{0.12\textwidth}|p{0.50\textwidth}|p{0.20\textwidth}|}
\hline
\textbf{Req} & \textbf{Elemento} & \textbf{Versión} \\ \hline
P1 & Clases & V1.0 \\ \hline
P2 & Actividades & V1.0 \\ \hline
P3 & Estados & V1.0 \\ \hline
P4 & Consistencia & V1.0 \\ \hline
P5 & Requisitos & V1.1 \\ \hline
P6 & MoSCoW & V1.0 \\ \hline
P7 & Validación & V1.0 \\ \hline
P8 & Pruebas & V1.1 \\ \hline
P9 & Trazabilidad & V1.1 \\ \hline
\end{tabular}

\vspace{0.5cm}

\textbf{Línea base:} Conjunto de documentos y modelos aprobados que se quedan congelados hasta que se apruebe otro cambio.


\end{document}
